\chapter[Sermon 11. The Latter Part of the Third Epistle Is Expounded, Wherein Is Spoken of the Nicolaitans, Which Are Damned. And Exhortation Is Made to Repentance.]{The Latter Part of the Third Epistle Is Expounded, Wherein Is Spoken of the Nicolaitans, Which Are Damned. And Exhortation Is Made to Repentance.}
\label{sermon:011}
\markright{Sermon 11. The Latter Part of the Third Epistle Is Expounded ...}

But I have a few things against you: that you hold there, that you maintain the doctrine of Balaam, which he taught in Balak, to put occasion of sin before the children of Israel, that they should eat of meat dedicated to idols, and commit fornication. Even so, I hate those who maintain the doctrine of the Nicolaitans. Repent, or else I will come to you shortly, and will fight against them with the sword of my mouth. Let him that has ears, hear what the Spirit says to congregations. To him that overcomes, I will give to eat manna that is hidden, and give him a white stone, and in the stone a new name write, which no man knows saving he that receives it.

\index{Augustine, Saint}In the first part of this letter, the Lord praises many things within the church of Pergamum, in the second part he will rebuke a few. And he says a few things, not that the error of the Nicolaitans is a minor offense, but that the sin is in others rather than in the true church itself: namely, in those who, though not truly members of the church’s body, outwardly joined with the church and wished to be considered members of it. After this, he speaks gently, lest by excessively exasperating the sin and error among the faithful, he should trouble their minds and discourage them utterly. There is a measure in all things, as the common saying is. And if in a church so commendable, something is found worthy of rebuke, what shall we say of those that are less commendable? Indeed, why should we not always see in all churches something to be blamed: not so much because the saints are always troubled with the weakness of the flesh, as because hypocrites and corrupt persons always join themselves to the church of God, such as were here the Nicolaitans, and as Judas the thief and traitor was among the apostles. In Christ, the church is without any spot or wrinkle, as the Lord says in John 13. And in the coming age it shall be most fully made perfect, which Saint Augustine also affirms.

And the Lord Jesus rebukes in the church of Pergamum, not that they uphold the Nicolaitan or Galaamitical doctrine, but that they harbor individuals who do. They were therefore culpable because they did not detest the Nicolaitans as much as the Ephesians did: of whom we heard in the first epistle, that they could not tolerate the wicked. Wherefore, lest the leaven of error should spread further through the entire mass of dough, the old leaven must be purged. It must be tested whether you favor or adhere to heresies. Furthermore, the Lord requires that we should not nurture them, but that we should pursue them with a holy hatred. Of this there is speaking in the first epistle.

Moreover, he describes the heresy of the Nicolaitans to the intent we may see, wherefore he blames it, wherefore he condemns it, and wherefore it ought to be hated. And he describes it carefully by the example of Scripture, that modest concerns or sensitivities might not be hurt or offended. I told you before how the ancient writers report the Nicolaitans to be most filthy. But all things are most aptly and carefully declared of Christ. They are taken from the 22nd, 23rd, 24th, and 25th chapters of the fourth book of Moses, called Numbers. He calls the Nicolaitan doctrine the doctrine of Balaam, and that by a comparison. In Balaam the soothsayer, these wicked acts are manifest; from which it may easily appear what sort of doctrine he had. First, he took the reward or price of iniquity, as St. Peter terms it, and would curse those whom God has blessed, doing clean contrary to his own mind. Secondly, he gives the king most pestilent counsel, which Scripture therefore calls a slander or offense. For he taught the king a way whereby he might entice the people of God into certain destruction, into the most unclean feeding of meats offered to idols, and into most filthy whoredom. All this shall be counted the doctrine of Balaam, which, in hope of filthy lucre and against God’s word and his own conscience, teaches idolatry, unclean eating, and fornication; or reproves not, but counsels rather, when he knows the thing to be filthy. Even so did the Nicolaitans, in speaking evil of the truth and of Christian purity, gave naughty counsel to many, that they should be partakers of meats offered up to idols, and couple with harlots, as in the first epistle I declared more at large.

Here we observe, by the example of our Savior Christ, how heresies should be refuted, not with angry or abusive words, but rather by passages and examples from holy Scripture. As is fitting here, the heresy of the Nicolaitans is condemned. And once condemned by the Lord, it remains condemned forever; we need no new deliberations to condemn impurity. Again, even if all the councils in the world decree the contrary, this remains true and certain, which the Lord Christ pronounces here: accursed be he who decides otherwise.

\index{Zurich}And here it seems good now to consider whether the Balaamitic and Nicolaitane doctrine in the church is completely extinguished. The names of Balaam and Nicolaitans we certainly abhor, but the thing itself, as well in the affairs of spiritual as temporal rulers, is openly present. For there are men in high authority, in various kinds of learning rightly excellent, most expert in the laws both of God and men, who nevertheless, blinded by the reward of iniquity, curse both the persons and things which they know that God has blessed. Of these Saint Peter also made mention in the second chapter of the second Epistle. The same do suggest evil counsels to Kings and Princes, tending to the destruction both of the preaching of the Gospel and the safeguard of the Church. The same being given to idolatry, and drowned in fleshly pleasures, eat of the sacrifices of the dead, and even feed of idol offerings, and in fornications run at riot. Consider, I pray you, what is the life and sustenance of most popish priests, what opinion they have of holy matrimony, and how much they abhor adultery and whoredom. They dare be bold to condemn matrimony and to judge whoredom better, so that they may enjoy the sacrifices of the dead and in many ways take their pleasure. If any, for the avoiding of whoredom, be joined in lawful matrimony, he is considered unworthy to sacrifice or to come at the altar; but whoremongers are admitted thick and threefold. And all they for the most part are the most beastly bondslaves of the belly, of whom you may believe that the holy Apostle of Christ, Saint Paul, has spoken: whose God is the belly, and glory in reproach of them that seek earthly things. And who will not acknowledge these and affirm them to be very Nicolaitans, maintaining the doctrine of Balaam the enchanter? Among the secular [? populace] you shall find men of all sorts who value the doctrine of Balaam and the wantonness of Zuryk more than they do modesty, gravity, and Christian sincerity. They love the liberty and wantonness of the flesh. They will not have youth and free people to be restrained by virtuous laws. They will even at this day banquet and mask with the maidens of Madian, and follow their fleshly lust. For they maintain surfeiting, drunkenness, and whoredom. And these are also very Nicolaitans, and have neither few nor obscure adherents to favor their sect. They want not their worldly reasons, many and great, to maintain the same.

But let us hear what Christ himself, sitting at the right hand of his Father, judges of them. Those or that same which these men think, teach, and do, I hate, says the Lord. What thing can be spoken more grievously than that God hates the doctrine of the Nicolaitans? For the whole scripture of both Testaments condemns this Nicolaitism.

\index{Repentance}After this description and reprehension of the Nicolaites, he proceeds, as in the former epistles, to exhort them to amendment or repentance. For when he says, “repent,” he understands or comprehends all penance or repentance. We said that to be a conversion unto God, whereby we amend evil things for good, relinquishing that which is evil and in its stead placing that which is good: and that through faith in sincere love and fear of God. You shall amend, therefore, if you abstain from meats offered up to idols and from fornication, and receive the true religion of Christ instituted, and possess your body in honor, not in the lust of cupidity, as Paul says in 1 Thessalonians 4. The church of Pergamum repeated, that if they did not disable or wink at the filthiness of the Nicolaites, but stoutly withstood it. The Nicolaites repented, if laying aside their filthiness, they received again the purity of faith and life. And to all and singular is said, “repent.”

The Lord now also drives them to repentance with severe threats: “Unless you amend,” he says, “I will come to you shortly,” a manner of speaking that has been discussed before. He adds, “And I will fight with you with the sword of my mouth.” With whom? With the impenitent, and especially with the Nicolaitans. He did not threaten utter destruction or devastation to the church, where there was great hope that the old leaven would be purged, but he threatens the impenitent people. Just as a judge, magistrate, or soldier uses the sword, so does Christ his word. And the word, indeed, wounds or slays no one, but it reveals God’s word, so does the execution of God’s power. Therefore, Christ, even as he reveals with his word, reveals that he will judge idolaters, those who worship false gods, swine, dogs, and adulterers, and not only judge them, but punish them. And as he threatens, he does.

Thus he fights with the sword of his mouth. We have an example in the Israelites, of whom 25,000 men were destroyed for following the doctrine of Balaam. Afterward, the Moabites and Midianites were also destroyed, nor were the corrupt women spared. Moses discusses this at length in Numbers 31. We also see at this day the sword of God going through the world, overthrowing these now, and then, for no other causes than those for which the Lord slew and destroyed Balaam and his adherents. Therefore, let us fear the Lord and walk in his commandments, for he will strike far more grievously with his sword when he pronounces judgment, “Go, you cursed, into everlasting fire.” Matthew 25. He does not expressly say, “I will cut you with the sword of my mouth.” For we are often severed and cut by the word of God to our great profit, discipline, and amendment. At present, he says he will fight: behold, he will fight, namely, against his enemies. Therefore, he threatens destruction. And we doubt nothing but that the impenitent of those times and of all times shall be destroyed, for (as I said even now) we lack no examples at this day.

Again, so that this notable and wholesome doctrine might not seem to belong only to a few men of Pergamum, he applies it to all churches throughout the world. We have discussed this application on one or two occasions in the previous letters.

Finally, in his manner, intending to move us all more strongly to repentance and obedience, he proposes a most ample promise: and that to those who strive and overcome the flesh, the world, and the devil, not to idlers, nor to those who wallow in wickedness. We are therefore encouraged by that promise, which is of three kinds. First, he promises to those who fight bravely and overcome, and do their duty, manna, and that secret or hidden [illegible]. That external manna, known to all men, is not the true manna. For the ungrateful Israelites despised it, considering it a very light food, and preferring the meat pots of Egypt, full of meat, onions, leeks, and garlic, that they might eat their fill. They do not see the celestial manna figured by this outward manna, yielding all sweetness and spiritual pleasure. The faithful see that this hidden manna is Christ, as he himself explains in John 6. Christ, therefore, gives himself to those who overcome, gives himself in food, which truly nourishes. He who has once, with true faith, tasted Christ will desire no other food to be given to him. For in Christ he has all things, in Christ he is complete, and with all good things fully satisfied. Oh, that our subtle disputers understood these things, they would reason about nothing at all concerning the merits and intercession of saints and such other things, about which, while they reason after their customary manner, they declare that they have not yet tasted how good and sweet the Lord is.

After he promises to give to the victors a white stone, that is, absolution and remission of all sins, and that without any doubt. For Christ truly absolves us from our sins and from the punishment due for them, and from condemnation. And he referred to the custom of ancient judges, who condemned with black stones and acquitted with white. For these verses are well known in Ovid’s \textit{Metamorphoses}, book 15. The practice long ago was to condemn with black stones and acquit with a white stone, and here we give warning that the remission of sins is not granted to people living for their works or merit, but that faith is the victory that overcomes the world. Which John himself testifies. And that faith indeed fights bravely in our hearts, but in the meantime acknowledges in all things the grace of God, neither does it nullify the merit of Christ. For as it is not sluggish, so is it also fearful.

Finally, he promises that he will write a new name on the stone, and that it is known only to him who possesses it. Christ will not only grant us remission of our sins, but also the glory and unspeakable joy of his heavenly communion. Isaiah and other prophets have made mention of this new name. Conquerors had famous names. If we overcome, we shall enjoy celestial glory. That glory is so immeasurable that it can only be perceived by feeling, not by speaking. For whatever you may say, however great, famous, or excellent, something greater will be given to the overcomers. For the Apostle Paul cites from Isaiah: What no eye has seen, nor ear heard, God has prepared for those who love him. And even in this present world, we are given a quiet conscience and unspeakable joy, which those who experience it truly feel. Those who have not tasted it cannot believe how great it truly is. Therefore, Paul says, “And the peace of God which surpasses all understanding,” etc. May our Savior Christ grant us such minds. Amen.
